% !TEX program = lualatex
% Auto-generated from YAML (v3) — 後期3: アルゴリズムと効率

\documentclass[aspectratio=43,professionalfonts,handout]{beamer}
\usepackage{beamer_template}

\newcommand{\LectureNum}{3}
\newcommand{\LectureTitle}{アルゴリズムと効率}
\newcommand{\CourseName}{情報リテラシー}
\newcommand{\TextPages}{pp.78-81}
\newcommand{\NextTopic}{情報通信ネットワーク（デジタル通信の仕組み）}
\newcommand{\NextPages}{対応なし}

\begin{document}
% --- Slide 1: title ---

\begin{frame}{\CourseName\quad 第\LectureNum 回「\LectureTitle 」}
  {\small \TermName\quad \TeacherName\quad ／\quad 教科書 \TextPages}

  \vfill

  \begin{block}{今日の流れ}
    \begin{enumerate}
      \item アルゴリズムとは
      \item 制御構造（順次・分岐・反復）
      \item フローチャートによる表現
      \item 探索アルゴリズム（線形・二分）
      \item 整列アルゴリズム
    \end{enumerate}
  \end{block}
\end{frame}

% --- Slide 2: terms ---

\begin{frame}{基本用語}
  \begin{description}[labelwidth=6em]
    \item[\kw{アルゴリズム（algorithm）}]
      ある問題を解決するために，一定の手順に従って解いていけるように表現されたもの
    \item[\kw{フローチャート（flow chart）}]
      処理手順を図記号を用いて図解したもの。流れ図ともいう
    \item[\kw{線形探索法}]
      先頭から順番に目的のデータを探す探索アルゴリズム
    \item[\kw{二分探索法}]
      探す範囲を半分にしながら目的のデータを探す探索アルゴリズム。データが昇順・降順に並んでいることが前提
    \item[\kw{選択ソート}]
      残りのデータから最小（最大）値を選び，順番に並べ替える整列アルゴリズム
  \end{description}
\end{frame}

% --- Slide 3: table ---

\begin{frame}{フローチャートの主な記号（JIS X 0121）}
  \centering
  \small
  \begin{tabular}{lll}
    \toprule
    記号（形） & 名称 & 内容 \\
    \midrule
    角丸四角形 & 端子 & 開始と終了 \\
    平行四辺形 & データ & データ入出力 \\
    四角形 & 処理 & 演算などの処理 \\
    台形型四角 & 表示 & 画面などに表示する \\
    ひし形 & 判断 & 条件による分岐 \\
    五角形/台形 & ループ端 & ループの始まり・終わり \\
    矢印線 & 線 & データや制御の流れ \\
    \bottomrule
  \end{tabular}
\end{frame}

% --- Slide 4: points ---

\begin{frame}{3つの制御構造}
  \begin{block}{ポイント}
    \begin{itemize}
      \item 順次構造：1つの処理が終わったら，次の処理へと順に実行する
      \item 分岐構造：条件によって処理を選択し，実行する
      \item 反復構造：条件によって同じ処理を繰り返し実行する（ループ）
    \end{itemize}
  \end{block}

  \vfill

  \begin{exampleblock}{補足}
    \begin{itemize}
      \item アルゴリズムはこの3つの制御構造と入出力の組み合わせで表現できる
    \end{itemize}
  \end{exampleblock}
\end{frame}

% --- Slide 5: twocol ---

\begin{frame}{探索アルゴリズムの比較}
  \begin{columns}[T]
    \begin{column}{0.48\textwidth}
      \begin{block}{線形探索法}
        \small
        \begin{itemize}
          \item 先頭から順番に1つずつ確認する
          \item データの並び順は問わない
          \item 最小：1回，最大：データ数と同じ回数
          \item 例：365日の誕生日当て → 最大365回
        \end{itemize}
      \end{block}
    \end{column}
    \begin{column}{0.48\textwidth}
      \begin{exampleblock}{二分探索法}
        \small
        \begin{itemize}
          \item 中央を確認し，大小で範囲を半分に絞る
          \item データが昇順（または降順）に並んでいることが前提
          \item 例：365日の誕生日当て → 最大9回で終わる
          \item データ数が多いほど効率の差が顕著になる
        \end{itemize}
      \end{exampleblock}
    \end{column}
  \end{columns}
\end{frame}

% --- Slide 6: points ---

\begin{frame}{整列アルゴリズム（選択ソート）}
  \begin{block}{ポイント}
    \begin{itemize}
      \item ① 全データの中から最小値を見つけ，先頭と入れ替える
      \item ② 残りのデータの中から最小値を見つけ，左から2番目と入れ替える
      \item ③ 以下，同様に繰り返す
      \item 比較・入れ替えの回数が少ないほど，より効率のよいアルゴリズムといえる
    \end{itemize}
  \end{block}

  \vfill

  \begin{exampleblock}{補足}
    \begin{itemize}
      \item 他の整列アルゴリズム：バブルソート（隣どうし比較），挿入ソート，マージソート，クイックソート
    \end{itemize}
  \end{exampleblock}
\end{frame}

% --- Slide 7: twocol ---

\begin{frame}{アルゴリズムの効率と「よいアルゴリズム」}
  \begin{columns}[T]
    \begin{column}{0.48\textwidth}
      \begin{block}{効率の評価}
        \small
        \begin{itemize}
          \item 同じ結果でも処理の回数が異なる
          \item データ数が増えたとき，手間がどれだけ増えるかで評価する
          \item 線形探索：データ数に比例（O(n)）
          \item 二分探索：データ数の対数に比例（O(log n)）
        \end{itemize}
      \end{block}
    \end{column}
    \begin{column}{0.48\textwidth}
      \begin{exampleblock}{よいアルゴリズムの条件}
        \small
        \begin{itemize}
          \item 正確である：目的通りの処理ができる
          \item 効率がよい：計算回数・メモリ使用量が少ない
          \item 読みやすい：人間にとって理解しやすい
          \item 再利用できる：他の問題にも応用できる
        \end{itemize}
      \end{exampleblock}
    \end{column}
  \end{columns}
\end{frame}

% --- Slide 8: exercise ---

\begin{frame}{演習}
  {\small 各自で取り組んでください。}

  \vfill

  \begin{block}{基本問題}
    \begin{itemize}
      \item 順次・分岐・反復の3つの制御構造を，身近な例を使って説明せよ
      \item 「自動販売機で商品を買う手順」をフローチャートで書け
      \item 5枚のカード（1,3,5,7,9）から数字5を探すとき，線形探索と二分探索それぞれ何回の比較が必要か
    \end{itemize}
  \end{block}

  \vfill

  \begin{exampleblock}{応用問題（余裕がある人）}
    \begin{itemize}
      \item 10枚のカードを選択ソートで昇順に並べ替える手順を書き，比較・入れ替えの回数を求めよ
    \end{itemize}
  \end{exampleblock}
\end{frame}

% --- Slide 9: assignment ---

\begin{frame}{課題・次回予告}
  \begin{importantbox}[課題]
    教科書 pp.78-81 を復習すること。フローチャートの記号を覚えておくこと
  \end{importantbox}

  \vfill

  \begin{block}{次回}
    「\NextTopic 」（教科書 \NextPages ）
  \end{block}

  \vfill

  {\small
  質問があれば授業後またはTeamsで受け付けます。}
\end{frame}

\end{document}
